\section{Algebra delle Matrici}

\begin{exercisebox}[title={Esercizio 1: Prodotto Righe per Colonne}]
Date le matrici:
$$ A = \begin{pmatrix} 1 & 2 & 0 \\ -1 & 3 & 1 \end{pmatrix}, \quad B = \begin{pmatrix} 2 & -1 \\ 0 & 1 \\ 3 & 4 \end{pmatrix} $$
Calcolare, se possibile, i prodotti $AB$ e $BA$.
\end{exercisebox}

\begin{solbox}
\textbf{Calcolo di $AB$:}
$A$ è $2 \times 3$, $B$ è $3 \times 2$. Il prodotto $AB$ sarà $2 \times 2$.
$$ AB = \begin{pmatrix} 1\cdot 2 + 2\cdot 0 + 0\cdot 3 & 1\cdot(-1) + 2\cdot 1 + 0\cdot 4 \\ -1\cdot 2 + 3\cdot 0 + 1\cdot 3 & -1\cdot(-1) + 3\cdot 1 + 1\cdot 4 \end{pmatrix} 
= \begin{pmatrix} 2 & 1 \\ 1 & 8 \end{pmatrix} $$

\textbf{Calcolo di $BA$:}
$B$ è $3 \times 2$, $A$ è $2 \times 3$. Il prodotto $BA$ sarà $3 \times 3$.
$$ BA = \begin{pmatrix} 2 & -1 \\ 0 & 1 \\ 3 & 4 \end{pmatrix} \begin{pmatrix} 1 & 2 & 0 \\ -1 & 3 & 1 \end{pmatrix} =
\begin{pmatrix} 
2(1)+(-1)(-1) & 2(2)+(-1)(3) & 2(0)+(-1)(1) \\
0(1)+1(-1) & 0(2)+1(3) & 0(0)+1(1) \\
3(1)+4(-1) & 3(2)+4(3) & 3(0)+4(1)
\end{pmatrix} $$
$$ = \begin{pmatrix} 3 & 1 & -1 \\ -1 & 3 & 1 \\ -1 & 18 & 4 \end{pmatrix} $$
\end{solbox}

\begin{exercisebox}[title={Esercizio 2: Inversa con Gauss-Jordan}]
Calcolare l'inversa della matrice:
$$ A = \begin{pmatrix} 1 & 0 & 1 \\ 1 & 1 & 2 \\ 0 & 1 & 2 \end{pmatrix} $$
\end{exercisebox}

\begin{solbox}
Affianchiamo la matrice identità e riduciamo:
$$ \left(\begin{array}{ccc|ccc} 
1 & 0 & 1 & 1 & 0 & 0 \\ 
1 & 1 & 2 & 0 & 1 & 0 \\ 
0 & 1 & 2 & 0 & 0 & 1 
\end{array}\right) $$
$R_2 \to R_2 - R_1$:
$$ \left(\begin{array}{ccc|ccc} 1 & 0 & 1 & 1 & 0 & 0 \\ 0 & 1 & 1 & -1 & 1 & 0 \\ 0 & 1 & 2 & 0 & 0 & 1 \end{array}\right) $$
$R_3 \to R_3 - R_2$:
$$ \left(\begin{array}{ccc|ccc} 1 & 0 & 1 & 1 & 0 & 0 \\ 0 & 1 & 1 & -1 & 1 & 0 \\ 0 & 0 & 1 & 1 & -1 & 1 \end{array}\right) $$
Ora risaliamo. $R_2 \to R_2 - R_3$:
$$ \left(\begin{array}{ccc|ccc} 1 & 0 & 1 & 1 & 0 & 0 \\ 0 & 1 & 0 & -2 & 2 & -1 \\ 0 & 0 & 1 & 1 & -1 & 1 \end{array}\right) $$
$R_1 \to R_1 - R_3$:
$$ \left(\begin{array}{ccc|ccc} 1 & 0 & 0 & 0 & 1 & -1 \\ 0 & 1 & 0 & -2 & 2 & -1 \\ 0 & 0 & 1 & 1 & -1 & 1 \end{array}\right) $$
Quindi $A^{-1} = \begin{pmatrix} 0 & 1 & -1 \\ -2 & 2 & -1 \\ 1 & -1 & 1 \end{pmatrix}$.
\end{solbox}

\begin{exercisebox}[title={Esercizio 3: Potenze di Matrici}]
Data $J = \begin{pmatrix} 0 & 1 \\ 0 & 0 \end{pmatrix}$, calcolare $J^2$ e $J^3$.
\end{exercisebox}

\begin{solbox}
$$ J^2 = \begin{pmatrix} 0 & 1 \\ 0 & 0 \end{pmatrix} \begin{pmatrix} 0 & 1 \\ 0 & 0 \end{pmatrix} = \begin{pmatrix} 0 & 0 \\ 0 & 0 \end{pmatrix} = \mathbf{0} $$
Di conseguenza, $J^3 = J^2 \cdot J = \mathbf{0} \cdot J = \mathbf{0}$.
La matrice $J$ è \textit{nilpotente} di indice 2.
\end{solbox}
